\documentclass[11pt]{article}
\usepackage[margin=1in]{geometry}
\usepackage[T1]{fontenc}
\usepackage[utf8]{inputenc}
\usepackage{lmodern}
\usepackage{microtype}
\usepackage{xurl}
\usepackage{setspace}
\usepackage{enumitem}
\usepackage{hyperref}
\setlength{\emergencystretch}{3em}
\onehalfspacing

\title{Referee Report for \textit{Econometrica}\\[4pt]
\large ``Necessity Entrepreneurship''}
\author{}
\date{March 2026}

\begin{document}
\maketitle

\section*{Recommendation}

Reject.

\medskip

The paper has a real and potentially important idea: entrepreneurship should not
be treated as a single extensive-margin object when weak labor-market outside
options can push households into low-scale business ownership. The distinction
between nonemployer self-employment and employer entrepreneurship is useful, and
the paper is clearly stronger now that it organizes the quantitative section
around that split.

That said, I do not think the paper is close to the publication threshold for
\textit{Econometrica} in its current form. The central contribution is still too
diffuse, the mapping from the structural object to observable data is too loose,
and the quantitative claims remain under-identified relative to the amount of
interpretation the paper wants to place on them. My view is therefore that the
paper is promising, but not publishable here in its current version.

\section*{Summary}

The paper develops a quantitative occupational-choice model with search frictions,
unemployment insurance, and two entrepreneurial margins: nonemployer
self-employment and employer entrepreneurship. The key conceptual distinction is
between ``opportunity'' entrepreneurs, who would choose entrepreneurship even
with a viable wage-work option, and ``necessity'' entrepreneurs, who choose it
only when that outside option is weak or unavailable. Quantitatively, lower
unemployment insurance raises entrepreneurship primarily through self-employment,
while removing worker unemployment risk reallocates activity toward employer
firms and raises output. The paper also includes a recession-style exercise based
on weaker labor-market prospects for low-education households.

The paper's best feature is that it insists entrepreneurship is a compositional
object. That is the right instinct and gives the project a more interesting
economic question than a simple extensive-margin business-ownership exercise.
However, the current manuscript does not yet convert that insight into a
sufficiently sharp and externally disciplined contribution for a general-interest
economics journal.

\section*{Major Comments}

\begin{enumerate}[leftmargin=*,label=\textbf{\arabic*.}]
\item \textbf{The paper still lacks one dominant question of general interest.}

The manuscript currently tries to do several things at once: distinguish
necessity from opportunity entrepreneurship, separate self-employment from
employer firms, study unemployment insurance, study labor-market risk, and say
something about recession dynamics. All of those are related, but the paper does
not yet force the reader to carry away one central question and one central
answer.

The strongest question seems to be something like: when weaker labor-market
outside options raise business ownership, does entrepreneurship become more or
less productive? If that is the real question, then the abstract, introduction,
and quantitative section should be rewritten around it much more aggressively.
At present the paper reads more like a broad research program than a tightly
targeted Econometrica paper.

\item \textbf{The empirical bridge to the structural object is too weak for the
claims the paper wants to make.}

The paper is right that necessity entrepreneurship is latent and not directly
observed. But once that is conceded, the paper needs a much sharper statement of
what empirical objects do and do not correspond to the model. The manuscript
currently leans on several margins:

\begin{itemize}[leftmargin=*]
\item nonemployer versus employer business counts,
\item high-propensity business applications,
\item GEM motive measures,
\item recession composition facts.
\end{itemize}

These are all informative, but they are noisy and conceptually distinct. The
paper does not yet establish a compelling observational bridge from these
objects to the structural necessity-opportunity distinction. In the current
version, the empirical section is best interpreted as motivation, not
measurement. That is acceptable for some papers, but then the manuscript must be
much more modest in how it describes the quantitative contribution.

For \textit{Econometrica}, I would want either a much tighter empirical
measurement argument, a clearer sufficient-statistics bridge, or a stronger
demonstration that the structural object delivers new disciplined implications
that are not already captured by reduced-form compositional margins.

\item \textbf{The quantitative interpretation is not yet disciplined enough.}

The manuscript moves among several distinct outcome objects:

\begin{itemize}[leftmargin=*]
\item entrepreneur share,
\item self-employed share among entrepreneurs,
\item average entrepreneur size,
\item output per household,
\item output per entrepreneur,
\item occasional language suggestive of efficiency or misallocation.
\end{itemize}

These objects are not interchangeable. The current draft is better than earlier
versions, but it still does not fully settle what the paper's welfare or
efficiency claim is. If the paper wants to talk about misallocation or
efficiency, then it needs either a welfare object or a more explicit allocation
criterion. If instead the contribution is about output and composition, the
paper should say that clearly and stop drifting into broader efficiency
language.

This matters because the headline result is subtle: lower unemployment insurance
raises entrepreneurship but reduces scale and output. That is interesting. But
the manuscript weakens the point by allowing the interpretation to slide between
``more necessity entry,'' ``more misallocation,'' and ``worse aggregate
outcomes'' without fully distinguishing them.

\item \textbf{The policy decomposition remains incomplete.}

The unemployment-insurance experiments still combine at least two channels:

\begin{itemize}[leftmargin=*]
\item the labor-market outside-option channel,
\item the fiscal-closure channel through taxes and transfers.
\end{itemize}

The draft is aware of this issue, but the interpretation still does not isolate
the mechanism cleanly enough. The fixed-tax comparison is therefore not a minor
robustness check. It is central to the paper's interpretation of the no-UI and
low-UI counterfactuals. Without a clean separation of these channels, the
reader cannot tell how much of the result reflects necessity entry per se and
how much reflects the fiscal consequences of changing the insurance system.

I would not be comfortable with publication in a journal of this level until the
paper treats this decomposition as part of the core evidence rather than as
appendix support.

\item \textbf{The new self-employment calibration is more plausible than before,
but the identification of the new parameters is still not persuasive enough.}

The revised benchmark is clearly better: closure risk, capital loss on failure,
and an employer-threshold cost give the model a more credible distinction
between self-employment and employer entrepreneurship. But the paper still does
not give the reader a sufficiently transparent parameter-to-moment mapping for
the key new objects. In particular, the manuscript needs a much clearer account
of how the following are pinned down:

\begin{itemize}[leftmargin=*]
\item entrepreneur closure risk by education,
\item capital loss on closure,
\item the employer-threshold fixed cost,
\item the self-employed versus employer composition target.
\end{itemize}

At the moment the calibration section reads as broadly reasonable, but not yet
fully convincing. For an Econometrica audience, it is not enough that the
choices feel sensible; the paper needs to distinguish much more clearly between
externally disciplined values, moment-matching parameters, and benchmark choices
used to close under-identified objects.

\item \textbf{The recession exercise is not yet publication-ready.}

The paper wants the recession section to show that weaker outside options shift
activity toward smaller-scale businesses. Conceptually that is a good extension.
But in the current manuscript, the recession object is still too unstable and
too weakly mapped to actual recession conditions.

Most importantly, the paper itself notes that a simple productivity MIT shock is
not the right object. That is correct. But the replacement recession exercise
also remains problematic, because the labor-market tightness response is not
cleanly recession-like, and the transition interpretation still appears to rely
on a partially ad hoc mapping from separation risk to recession conditions.

In other words, the current recession section does not yet read like a decisive
quantitative result. It reads like an extension that is still under
construction. For this reason, I would either sharply narrow the claim or remove
the section from the main text until the recession mapping is much more
convincing.

\item \textbf{The paper remains too cluttered for a top general-interest
submission.}

The manuscript is strongest when it focuses on four disciplined margins:

\begin{itemize}[leftmargin=*]
\item total entrepreneurship,
\item self-employed share among entrepreneurs,
\item average entrepreneurial scale,
\item aggregate output.
\end{itemize}

It is weaker when the main text becomes crowded with large by-education tables,
cutoff details, and multiple partially overlapping quantitative objects. Those
materials are useful for development and for appendices, but the current draft
still feels like it is trying to show the reader every supporting object at
once. That is not how papers at this level are usually packaged.

The main text needs much more selectivity. Right now the manuscript often reads
as though it is proving seriousness by showing more objects, when in fact it
would be stronger if it showed fewer objects and interpreted them more sharply.
\end{enumerate}

\section*{Minor Comments}

\begin{enumerate}[leftmargin=*,label=\textbf{\arabic*.}]
\item The paper's external presentation still needs substantial polishing.
Although the draft is much cleaner than before, the manuscript does not yet read
like a finished top-journal submission. The prose is often serviceable rather
than sharp, some sections still have note-like transitions, and the paper would
benefit from a substantial editorial pass.

\item The calibration table and related text should state more explicitly which
rows are benchmark parameter values and which are equilibrium outcomes. This is
especially important for taxes and other fiscal objects.

\item The manuscript should be more explicit about what the benchmark business
type split is intended to match in the data. If the paper is going to
distinguish self-employed nonemployers from employer firms, the empirical target
for that split should be stated clearly and repeatedly.

\item The paper should decide whether it is primarily about occupational choice,
entrepreneurial composition, or insurance design. My read is that the best
version is a paper about entrepreneurial composition under changing outside
options. That focus should discipline both the introduction and the conclusion.
\end{enumerate}

\section*{Overall Assessment}

This is not a weak paper. It contains a legitimate idea, and the distinction
between lower-scale self-employment and employer entrepreneurship is a genuine
improvement over more aggregated entrepreneurship exercises. I can also see a
path by which the project becomes a strong paper.

But I do not think that path leads to \textit{Econometrica} in its current form.
The paper would need a much sharper statement of the contribution, a much more
disciplined empirical bridge, cleaner quantitative identification, and a more
convincing treatment of the recession and policy decompositions. My bottom-line
assessment is therefore:

\begin{itemize}[leftmargin=*]
\item promising project,
\item materially improved relative to earlier drafts,
\item not publishable at \textit{Econometrica} in its current state.
\end{itemize}

My guess is that with substantial refocusing and a cleaner empirical bridge,
this project could become a good field-journal paper and possibly stronger than
that if the authors can convert the necessity-opportunity distinction into a
much more sharply testable quantitative claim.

\section*{Most Important Revisions}

If the authors continue developing the paper, I would prioritize the following:

\begin{enumerate}[leftmargin=*]
\item Rewrite the abstract and opening introduction around one question and one
answer.
\item Tighten the empirical bridge and explicitly distinguish measurement from
motivation.
\item Make the fixed-tax versus endogenous-tax decomposition central to the
policy interpretation.
\item Simplify the main text and move more quantitative detail to appendices.
\item Clarify the calibration mapping from moments to the new self-employment
parameters.
\item Either sharpen the recession exercise substantially or demote it to a more
limited robustness role.
\end{enumerate}

\end{document}
